% Review draft: FF.45. Not separately input after integration.
\subsection{Faithfulness of generic restriction on the pointwise kernel}

\begin{claim}{FF.45}{Proved}\label{clm:FFgenerickernel}
Let \(k\) be a characteristic-zero field, let \(X/k\)
be a smooth geometrically integral curve, and put
\(F=k(X)\). A pinned automorphism
\(\varphi:X_{\mathrm{\acute et}}\to
X_{\mathrm{\acute et}}\) over \(BG_k\) induces a generic
outer automorphism
\[
 \alpha_\varphi\in
 \operatorname{Aut}_{G_k}(G_F)/
              \operatorname{Inn}(G_{\bar k(X)}).
                                                        \tag{FF.45.1}
\]
If its actual function on \(X(\bar k)\) is the identity
and (FF.45.1) is trivial, then \(\varphi\) is isomorphic
to the identity over \(BG_k\).

Consequently, for every finite \(E/B_p\), the remaining
pointwise kernel in Claim~\ref{clm:FFmetabelianaffine}
injects into the relative birational outer automorphism
group of \(G_{E(T)}\). Its image preserves the geometric
conjugacy class of every finite-point inertia subgroup,
and its induced map on the corresponding residue Galois
group is the identity. In particular, an inner generic
action cannot hide a nontrivial automorphism of the full
etale topos.
\end{claim}

\begin{proof}
Write \(j:\operatorname{Spec}F\to X\), and choose a
separable closure of \(F\) containing \(\bar k\).
We retain all structural isomorphisms over \(BG_k\).

\emph{The generic action.}
An automorphism preserves the unique generic point and
therefore permutes the nonempty open subtopoi. The filtered
category of quasicompact separated etale objects over
these opens, with objects and arrows identified after
further shrinking, is the category of finite etale
\(F\)-schemes. Indeed, every finite separable extension
of \(F\), and every morphism between finitely many such
extensions, spreads out over a nonempty open of \(X\).
Conversely, two morphisms that agree generically agree
after shrinking. Thus \(\varphi\) induces an
autoequivalence of finite continuous \(G_F\)-sets.
The specified relative structure fixes the subcategory
of sets inflated from \(G_k\), with its specified
comparison. Galois-category reconstruction gives exactly
(FF.45.1), with inner ambiguity only in the geometric
kernel.

Assume from now on that the geometric point function is
the identity and (FF.45.1) is trivial. Choose the resulting
natural identity for the generic inverse-image functor,
over the category of \(G_k\)-sets. It extends from finite
continuous \(G_F\)-sets to all continuous \(G_F\)-sets:
each such set is a union of its finite orbits.

\emph{The maximal etale extension of a generic object.}
For a finite continuous \(G_F\)-set \(S\), the sheaf
\(j_*S\) is represented by the etale locus of the finite
normalization of \(X\) in the finite etale \(F\)-algebra
corresponding to \(S\). The normalization is finite
because \(X\) is a finite-type regular curve over a
field. To check the assertion, an \(F\)-map from the
generic fibre of an etale normal \(X\)-scheme extends
uniquely to its normalization map. At every point in its
image the extension of discrete valuation rings has
ramification index one and separable residue extension,
so the image lies in the etale locus. This gives precisely
the defining sections of \(j_*S\).

Every quasicompact separated etale \(X\)-scheme \(U\)
with generic fibre \(S\) is an open subobject
\[
                           U\ \lhook\joinrel\longrightarrow j_*S.
                                                        \tag{FF.45.2}
\]
For example, this follows by applying the same
normalization argument to each component: a separated
quasi-finite etale map into a normal curve is an open
subscheme of its finite normalization. All components of
\(U\) dominate \(X\), since etale maps are open.
The complement in the etale normalization locus consists
of finitely many closed points.

For a geometric point \(\bar x\) above a closed point
\(x\), choose a prolongation of its valuation to the
separable closure of \(F\). Let \(D_x\) and \(I_x\)
be the corresponding decomposition and inertia groups.
Then
\[
               (j_*S)_{\bar x}=S^{I_x},\qquad
               D_x/I_x=G_{\kappa(x)}.
                                                        \tag{FF.45.3}
\]
The stalk \(U_{\bar x}\) in (FF.45.2) is a
\(D_x\)-stable subset of \(S^{I_x}\). This description
also follows directly by base change to the strict
henselization at \(\bar x\): its unramified sheets are
exactly the \(I_x\)-fixed generic sheets.

\emph{The relevant normalizer.}
In this setting
\[
                            N_{G_F}(I_x)=D_x.
                                                        \tag{FF.45.4}
\]
Indeed, the fixed field \((F^{\rm sep})^{I_x}\) is the
fraction field of the strict henselization at the chosen
place. It has a henselian discrete rank-one valuation,
separably closed residue field of characteristic zero,
and value group \(\mathbb Z\). It is not separably
closed: a uniformizer has no square root in it. An element
normalizing \(I_x\) transports this valuation to another
rank-one henselian valuation on the same field. If the
two valuations were distinct, they would be incomparable,
and the two-valuation argument proved in
Claim~\ref{clm:FFstricttest} would make the field separably
closed. They therefore coincide. Its unique prolongation
to \(F^{\rm sep}\) is then preserved as well, so the
normalizing element belongs to \(D_x\). The reverse
inclusion is the usual normality of inertia in the
decomposition group.

\emph{The specialization comparison is forced.}
The chosen generic identity and the adjunction
\(j^*\dashv j_*\), together with invertibility of
\(\varphi^*\), give natural identifications
\[
                             \varphi^*j_*S\simeq j_*S.
                                                        \tag{FF.45.5}
\]
Actual pointwise identity gives an isomorphism
\(\bar x^*\varphi^*\simeq\bar x^*\), by
\citet[Proposition~3.6]{FF_CHW2024}. Combining it with
(FF.45.5) yields a natural automorphism of the functor
\[
       \{\text{finite continuous }G_F\text{-sets}\}
          \longrightarrow\{\text{sets}\},\qquad
                            S\longmapsto S^{I_x}.
                                                        \tag{FF.45.6}
\]
There is no arbitrary specialization permutation here.
The functor (FF.45.6) is prorepresented by the profinite
\(G_F\)-set \(G_F/I_x\). Concretely, write it as the
inverse system of finite \(G_F\)-sets
\(G_F/(I_xN)\), where \(N\) ranges over open normal
subgroups. A continuous map from this system to a finite
\(G_F\)-set is exactly an \(I_x\)-fixed element.
Pro-Yoneda therefore identifies the natural automorphisms
of (FF.45.6) with
\[
                      N_{G_F}(I_x)/I_x=D_x/I_x.
                                                        \tag{FF.45.7}
\]
Equivalently, the comparison acts on every \(S^{I_x}\)
by one element of \(D_x/I_x\). Hence it preserves
every \(D_x\)-stable subset, in particular
\(U_{\bar x}\). If the relative point comparisons are
chosen explicitly over \(BG_k\), their action on sets
inflated from \(G_k\) is the identity; this even makes
the element in (FF.45.7) trivial, since
\(G_{\kappa(x)}\to G_k\) is injective. Preservation of
the subsets is all that is needed.

It follows from (FF.45.5)--(FF.45.7) that \(\varphi^*\)
fixes the subobject (FF.45.2) inside \(j_*S\): its
generic stalk is fixed by construction and all its closed
geometric stalks are fixed by the preceding paragraph.
Geometric points detect equality of subobjects. This
produces an isomorphism \(\varphi^*U\simeq U\) for
every quasicompact separated etale \(U\). These
isomorphisms are natural. A morphism between two such
objects is determined by its generic restriction, because
the source is reduced with all components dominant and
the target is separated. Both ways around the proposed
naturality square agree generically and hence everywhere.

The quasicompact separated etale objects generate the
etale topos, and automorphisms preserve this subcategory;
see \citet[Theorem~A.8, Corollaries~A.9 and~A.11]
{FF_CHW2024}. Thus the natural identity extends to the
whole inverse-image functor. It respects the specified
structure over \(BG_k\): for a finite \(G_k\)-set the
corresponding constant etale cover is already its own
maximal extension \(j_*S\), and the generic identity
was chosen with precisely that relative compatibility.
This proves the first assertion.

For the final assertion take \(X=\mathbb A^1_E\).
The same specialization calculation, without assuming
the generic action trivial, identifies the prorepresented
functors attached to \(I_x\) and its transform. Thus
the generic action preserves the inertia conjugacy class
at each fixed geometric point. The relative point
comparison makes the conjugator geometric and makes the
map on the residue Galois group the identity. The first
assertion proves injectivity on the pointwise kernel.
\end{proof}

\paragraph{Exact birational scope at the metabelian field.}
For every finite \(E/B_p\), Koenigsmann's
\citet[Theorem~4.1, Corollaries~4.3 and~4.5]
{FF_Koenigsmann2003} apply: \(E\) is neither separably
closed nor real closed, and has a rank-one mixed
characteristic \((0,p)\) henselian valuation with
\(p\) dividing the order of \(G_E\).
They characterize the geometric decomposition and inertia
groups of the complete curve and the fundamental-group
quotients of its opens. Therefore the generic action in
Claim~\ref{clm:FFgenerickernel}, already fixing every
finite geometric place, fixes the remaining place at
infinity as well. Corollary~4.5 deduces a birational
Isom conclusion only when a separate fundamental-group
Isom conclusion is available for an open. It does not
prove that the pointwise birational stabilizer here is
trivial.

The more recent relative and semi-absolute theorems of
\citet[Theorems~A and~E, Definition~1.1,
Proposition~2.6]{FF_Tsujimura2026} retain an abelian-variety
Kummer-faithfulness hypothesis. For \(\Sigma=\{p\}\)
that hypothesis fails here: the elliptic curve in
Claim~\ref{clm:FFmetabelianbarrier} has trivial
\(G_{B_p}\)-action on its full \(p\)-adic Tate module,
so its nonzero rational \(p\)-power torsion points have
compatible \(p\)-divisions in \(B_p\). They are not
prime-to-\(p\) torsion. In fact every nonempty
\(\Sigma\) is excluded. If \(\ell\in\Sigma\)
and \(\ell\ne p\), the pro-\(p\) image of
\(G_{B_p}\) on the Tate module of any positive-dimensional
abelian variety injects into \(\operatorname{GL}_n(\mathbb F_\ell)\):
its intersection with the pro-\(\ell\) principal
congruence kernel is trivial. The image is therefore
finite. After a finite base extension all
\(\ell\)-power torsion is rational. A nonzero point
of this divisible \(\ell\)-primary subgroup is
\(\Sigma\)-divisible, since multiplication by each
other prime is bijective on that subgroup, but is not
\(\Sigma'\)-torsion. For \(\Sigma=\{p\}\) use
the preceding Tate example. This is a failure of that
source's hypothesis; the positive fields already proved
above show why it is not a reconstruction counterexample.
No Isom theorem for \(B_p\),
or vanishing of the remaining stabilizer, is inferred
from either source.

The commensurator theorem of
\citet[Theorem~6.8]{FF_HoshiMochizukiTsujimura2025}
also assumes TKND-AVKF. Its more general
Theorem~5.15 assumes QAA; Definition~5.12(ii)--(iii)
retains conditions on proper-cover cohomology and
function fields on all covers. These are not supplied
by the genus-zero conclusions below.
